\documentclass[12pt]{article}

\usepackage[utf8]{inputenc}
\usepackage[russian]{babel}
\usepackage{amsmath}
\usepackage{amsfonts}
\usepackage{amssymb}

\title{Пример документа LaTeX}
\author{Автор}
\date{\today}

\begin{document}

\maketitle

\begin{abstract}
Данная статья представляет собой пример документа LaTeX с введением и заключением. В ней демонстрируется структура стандартного научного текста.
\end{abstract}

\section{Введение}
Введение представляет собой начальный раздел работы, в котором формулируется тема, цель и актуальность исследования. В этой части читатель знакомится с основными понятиями и контекстом рассматриваемой проблемы. Введение должно заинтересовать читателя и обозначить рамки предстоящего анализа.

В современном научном тексте введение играет важную роль, так как именно с него начинается восприятие всей работы. В этой части автор должен четко сформулировать задачи исследования, обозначить методы, которые будут использованы, и кратко описать структуру статьи.

\section{Основная часть}
Основная часть статьи содержит детальный анализ рассматриваемой темы. Здесь приводятся теоретические сведения, экспериментальные данные, математические выкладки и другие элементы, зависящие от специфики исследования. В рамках этой работы основная часть может быть разделена на несколько подразделов, каждый из которых посвящен определенному аспекту темы.

\section{Заключение}
Заключение является завершающей частью работы, в которой подводятся итоги исследования. В этой секции автор формулирует основные выводы, обобщает полученные результаты и указывает возможные направления дальнейших исследований. Заключение должно логично вытекать из основной части и подтверждать достижение целей, поставленных во введении.

В заключение можно также отметить практическую значимость проделанной работы и предложить рекомендации по использованию полученных результатов в дальнейших научных или прикладных исследованиях.

\end{document}